% from ega4-10-fr.tex
\oldpage[IV-3]{103}
cet ensemble est \emph{très dense} dans $X$, donc $X$ est un espace de
Jacobson par définition.

Rappelons (5.1.10) que l'hypothèse faite sur $X$ dans (10.4.5.1) est toujours
vérifiée lorsque $X$ est l'espace sous-jacent à un \emph{préschéma}.

Revenant alors à la démonstration de (10.4.5), la condition \emph{b)} entraîne
que pour tout point $x$ non fermé de $\operatorname{Spec}(A)$, le point
générique $x$ de
$\operatorname{Spec}(A/\mathfrak j_x)=V(\mathfrak j_x)=\overline{\{x\}}$
n'est pas isolé dans $\overline{\{x\}}$, donc \emph{b)} entraîne \emph{a)} en
vertu du lemme (10.4.5.1) et de (10.4.4).

\begin{corollary}[10.4.6]
\label{IV.10.4.6-fr}
Toute algèbre de type fini $B$ sur un anneau de Jacobson $A$ est un anneau de
Jacobson, et l'image réciproque dans $A$ de tout idéal maximal de $B$ est un
idéal maximal de $A$. En particulier, toute algèbre de type fini sur un corps
ou sur $\mathbf Z$ est un anneau de Jacobson.
\end{corollary}

Une $B$-algèbre de type fini $K$ qui est un corps est aussi une $A$-algèbre de
type fini, donc un $A$-module de type fini et \emph{a fortiori} un $B$-module
de type fini, d'où la première assertion~; la seconde a été démontrée au cours
de la démonstration de (10.4.5), appliqué à $K=B/\mathfrak m$.

\begin{corollary}[10.4.7]
\label{IV.10.4.7-fr}
Si $X$ est un préschéma de Jacobson, $f:Y\to X$ un morphisme localement de type
fini, alors $Y$ est un préschéma de Jacobson et l'image par $f$ de tout point
fermé dans $Y$ est un point fermé dans $X$.
\end{corollary}

La question étant locale sur $X$ et sur $Y$, on est ramené au cas où $X$ et
$Y$ sont affines, et le corollaire résulte alors de (10.4.6).

\begin{corollary}[10.4.8]
\label{IV.10.4.8-fr}
Si $k$ est un corps algébriquement clos, $X$ un $k$-préschéma localement de
type fini, l'ensemble des points de $X$ rationnels sur $k$ est très dense dans
$X$.
\end{corollary}

En effet, $X$ est un préschéma de Jacobson (10.4.7) et les points fermés de
$X$ sont exactement les points rationnels sur $k$ (\textbf{I}, 6.4.2).

\begin{env}[10.4.9]
\label{IV.10.4.9-fr}
Le fait que les préschémas localement de type fini sur un corps ou sur
$\mathbf Z$ sont des préschémas de Jacobson est particulièrement important,
en raison de la possibilité de se ramener à ce cas dans de nombreuses
questions de Géométrie algébrique (8.1.2, \emph{c)}). Nous allons en donner
deux exemples~:
\end{env}

\begin{env}[10.4.10]
\label{IV.10.4.10-fr}
\textbf{Applications (10.4.10)~: I~: Démonstration de (6.15.9).} --- Soient
$k$ un corps séparablement clos, $X$ un $k$-préschéma localement de type fini
sur $k$ et \emph{unibranche}. On sait que la fermeture intégrale d'une
$k$-algèbre intègre de type fini $A$ dans une extension finie de son corps des
fractions est une $A$-algèbre finie (Bourbaki, \emph{Alg. comm.}, chap.~V,
\S~3, n\textsuperscript{o}~2, th.~2), donc toute $k$-algèbre de type fini est
un anneau universellement japonais~; on en conclut que l'ensemble des points
$x\in X$ où $X$ est géométriquement unibranche est localement constructible
(9.7.10). Mais l'hypothèse et le lemme (6.15.8) entraînent que cet ensemble
contient tous les points \emph{fermés} de $X$. La conclusion résulte donc de
(10.4.6), (10.3.1) et de la bijectivité de l'application canonique
$\mathfrak{Lc}(X)\to\mathfrak{Lc}(X_0)$ (où $X_0$ est l'ensemble des points
fermés de $X$) (10.2.7).
\end{env}

\begin{env}[10.4.11]
\EGAReaderTarget{ega4.c04.statement.10-4-11}{IV.10.4.11-fr}
\textbf{Applications~: II~: Proposition (10.4.11).} --- Soient $S$ un
préschéma, $X$ un $S$-préschéma de type fini. Tout $S$-endomorphisme de $X$
qui est radiciel est surjectif (donc bijectif).
\end{env}

